\documentclass[12pt]{article}
\usepackage[utf8]{inputenc}
\usepackage[T1]{fontenc}
\usepackage[margin=1in]{geometry}
\usepackage{setspace}
\usepackage{xcolor}
\usepackage[colorlinks=true, urlcolor=blue!60!black, linkcolor=blue!60!black]{hyperref}
\setlength{\parindent}{2em}
\setlength{\parskip}{1.25em}
\renewcommand{\baselinestretch}{1.0}
\begin{document}
\begin{center}
{\LARGE \textit{Missouri v. Holland}, 252 U.S. 416 (1920)\par}
\vspace{0.8em}
{\large Case Brief --- Four-Page Edition\par}
\end{center}
\vspace{1em}

\section{Memory Jogger}

Upholding the Migratory Bird Treaty Act 7--2, Justice Holmes held that Congress may enact legislation necessary and proper to execute a valid treaty and that Missouri's Tenth Amendment objection failed even assuming Congress could not regulate the birds without the treaty (\href{https://supreme.justia.com/cases/federal/us/252/416/}{Justia}).

\section{Detailed Case Facts}

In 1913 Congress tried to regulate migratory bird hunting without any treaty, but two district courts struck the law down as beyond Congress's powers, reasoning under \textit{Geer v. Connecticut} that wild game belonged to the States. The United States and Great Britain, acting for Canada, then concluded the Migratory Bird Treaty of 1916, which recited that many species migrating between the two countries were valuable as food and as destroyers of harmful insects and were in danger of extermination. It set closed seasons and bound both nations to propose implementing measures. Congress passed the Migratory Bird Treaty Act of 1918, barring the killing, capture, or sale of covered birds except as the Secretary of Agriculture's regulations allowed.

\section{Procedural History}

Missouri sued in equity in the Western District of Missouri to enjoin U.S. Game Warden Ray P. Holland from enforcing the Act, arguing that it invaded powers reserved by the Tenth Amendment and the State's title to wild birds. The district court dismissed, holding the Act constitutional, and Missouri appealed directly. The Supreme Court heard argument March 2, 1920, and affirmed on April 19, 1920.

\section{Judicial Votes}

Affirmed, \textbf{7--2}.
\begin{itemize}
\item \textbf{Majority (7):} Justice Holmes, joined by Chief Justice White and Justices McKenna, Day, McReynolds, Brandeis, and Clarke.
\item \textbf{Dissents:} Justices Van Devanter and Pitney, without opinion.
\end{itemize}

\section{Holding}

Missouri's quasi-sovereign interest in regulating game supported equity jurisdiction. The treaty was a valid exercise of the Article~II treaty power, and the Act was valid under Article~I, Section~8, as a necessary and proper means of executing it. Because the treaty power is expressly delegated and treaties are supreme law under Article~VI, a valid treaty and its implementing legislation are not limited to what an unaided Act of Congress could do and do not violate the Tenth Amendment.

\section{Analysis of Opinions}

\textbf{Majority (Holmes).} The Tenth Amendment reserves only powers ``not delegated to the United States,'' and the treaty power is delegated. Article~VI makes Acts of Congress supreme only when made in pursuance of the Constitution, but treaties when made under the authority of the United States; Holmes disclaimed an unlimited treaty power but said its limits must be found ``in a different way.'' The Constitution is ``a constituent act'' that ``called into life a being the development of which could not have been foreseen completely by the most gifted of its begetters.'' Since the treaty breached no ``prohibitory words,'' the only question was whether ``some invisible radiation'' from the Tenth Amendment forbade it. A problem of ``the sharpest exigency for the national wellbeing'' that States cannot solve alone requires national power. Missouri's claim of title failed because wild birds are possessed by no one, and treaties had long overridden state law (\textit{Ware v. Hylton}; \textit{Geofroy v. Riggs}). Relying on the States was ``vain.''

\textbf{Van Devanter and Pitney (dissenting).} No opinion. Later scholarship and the \textit{Bond} concurrences voiced the structural objection that treaties should not enlarge Congress's domestic powers.

\textbf{Later developments.} \textit{Reid v. Covert} (1957) held that no agreement with a foreign nation can free any branch from the Constitution's restraints. In \textit{Bond v. United States} (2014), the Court construed the chemical-weapons statute not to reach ``purely local crimes'' and did not revisit \textit{Holland}; Justice Scalia called its necessary-and-proper bridge ``unreasoned and citation-less,'' Justice Thomas would limit treaties to ``matters of international intercourse,'' and Justice Alito to ``matters of legitimate international concern.''

\section{Examples: Future Applications}

\begin{enumerate}
\item \textbf{Same-Side: Migrating Butterfly.} The United States, Canada, and Mexico agree to protect an endangered butterfly's breeding habitat, and Congress forbids habitat destruction on private land. As with the birds, the creature is in international transit and no State can protect it alone.
\item \textbf{Same-Side: Subway Nerve-Agent Attack.} A defendant releases nerve agent in a subway. That is core Convention conduct, \textit{Bond}'s narrowing does not apply, and \textit{Holland} supplies the authority.
\item \textbf{Opposite-Side: Treaty Banning Religious Insult.} A treaty requires criminalizing speech insulting religion. The First Amendment is a ``prohibitory'' limit, and under \textit{Reid} the statute is invalid as applied to protected speech.
\item \textbf{Opposite-Side: Poisoned Fish Pond.} A homeowner pours pesticide into a neighbor's pond. As in \textit{Bond}, the statute is read not to reach purely local crime, so the constitutional question never arises.
\item \textbf{Fence-Sitter: Treaty Ban on School Corporal Punishment.} After a ban on school corporal punishment fails under \textit{Lopez}, a ratified treaty requires one and Congress reenacts it with a clear statement. \textit{Holland} supports it, but the \textit{Bond} concurrences' limits---international concern and Scalia's reading of the Necessary and Proper Clause---cut against it.
\end{enumerate}

\section{Critique}

\textbf{Textual slippage.} Nicholas Quinn Rosenkranz, cited by Justice Scalia, argues that ``under the Authority of the United States'' was drafted to preserve treaties made under the Articles of Confederation, not to exempt future treaties from constitutional limits, and that Holmes never explained why the Necessary and Proper Clause reaches a treaty's substantive obligations rather than the process of making treaties.

\textbf{Functionalism without a stopping point.} Holmes's collective-action rationale fits birds crossing the U.S.--Canada border, but his opinion offers no manageable standard for separating genuine international problems from pretextual treaties designed to evade Article~I.

\section{Key Quotations}

\begin{quote}
``It was enough for them to realize or to hope that they had created an organism; it has taken a century and has cost their successors much sweat and blood to prove that they created a nation. The case before us must be considered in the light of our whole experience, and not merely in that of what was said a hundred years ago.''
\end{quote}
\noindent --- Justice Holmes, 252 U.S. at 433.

\begin{quote}
``If the treaty is valid, there can be no dispute about the validity of the statute under Article~I, \S~8, as a necessary and proper means to execute the powers of the Government.''
\end{quote}
\noindent --- Justice Holmes, 252 U.S. at 432.

\begin{quote}
``To put the claim of the State upon title is to lean upon a slender reed. Wild birds are not in the possession of anyone, and possession is the beginning of ownership. The whole foundation of the State's rights is the presence within their jurisdiction of birds that yesterday had not arrived, tomorrow may be in another State, and, in a week, a thousand miles away.''
\end{quote}
\noindent --- Justice Holmes, 252 U.S. at 434.

\end{document}
